\documentclass[10pt,a4paper,sans]{moderncv}
\moderncvstyle{banking}
\moderncvcolor{black}
\nopagenumbers{}

\usepackage[utf8]{inputenc}
\usepackage[scale=0.9, top=0.7cm, bottom=0.45cm]{geometry}
\usepackage{enumitem}
\usepackage{ragged2e}
\usepackage{xcolor}
\usepackage[unicode]{hyperref}

\hypersetup{
    colorlinks=true,
    urlcolor=blue
}

\name{Arjun}{Juneja}

\newcommand*{\customcventry}[7][.05em]{
\begin{tabular}{@{}l}
{\bfseries #4} \\
{\itshape #3}
\end{tabular}
\hfill
\begin{tabular}{l@{}}
{\bfseries #5} \\
{\itshape #2}
\end{tabular}
\ifx&#7&%
\else{\\
\begin{minipage}{\maincolumnwidth}%
\small#7%
\end{minipage}}\fi%
\par\addvspace{#1}}

\setlist[itemize]{leftmargin=0.55cm, itemsep=0.15mm, topsep=0.2mm}

\begin{document}
\vspace{1mm}

\begin{center}
    {\huge \textbf{Arjun Juneja}} \\[1mm]
    London, UK \textbar{}
+44 7447 361 102 \textbar{}
\href{mailto:arjunjuneja4u@gmail.com}{arjunjuneja4u@gmail.com} \textbar{}
\href{https://arjunjuneja.com}{arjunjuneja.com} \textbar{}
\href{https://linkedin.com/in/arjunjuneja}{LinkedIn} \textbar{}
\href{https://github.com/ImArjunJ}{GitHub}
\end{center}

\vspace{-6mm}

\section{Summary}
\justifying{
Systems and security engineer building low-level tooling across Windows kernel, Linux drivers, reverse engineering, and C++/Rust infrastructure. Published kernel-assisted threat detection research on arXiv, built open-source CAN tooling with 80+ GitHub stars, and won multiple international hackathon/CTF competitions including Junction '25 and LauzHack '25.
}

\section{Experience}

\customcventry{Jul 2026 -- Sep 2026}{Security Engineer Intern}{XTX Markets}{London, UK}{}{
\begin{itemize}
    \item Selected for Summer 2026 Security Engineering internship focused on security tooling and infrastructure.
\end{itemize}}

\vspace{0.4mm}

\customcventry{Sep 2025 -- Jan 2026}{Software Engineer}{Hacktron AI}{London, UK}{}{
\begin{itemize}
    \item Built tree-sitter-based static call graph extraction engine to identify security-relevant call paths for automated pentesting reports.
    \item Shipped CLI authentication flow supporting OAuth2 token exchange, persistent sessions, and developer-facing access control.
\end{itemize}}

\vspace{0.4mm}

\customcventry{Jul 2023 -- Sep 2023}{Software Engineer Intern}{BP}{London, UK}{}{
\begin{itemize}
    \item Built internal Safety \& Assurance tooling adopted by 100+ assurers across the organisation.
    \item Internship extended by 2 months after delivering production tooling and stakeholder-facing features.
\end{itemize}}

\section{Selected Projects \& Publications}
\cventry{}{jcan}{}{}{}{
Open-source CAN bus diagnostic tool with 80+ GitHub stars. Implemented Linux support for Vector CAN hardware with real-time bus monitoring, DBC decoding, signal plotting, and multi-adapter support. Built in C++23 with Dear ImGui and libusb.
\href{https://github.com/ImArjunJ/jcan}{Project}
}


\vspace{0.7mm}


\cventry{}{RX-INT: Kernel-Assisted Threat Detection}{}{}{}{
Built a Windows kernel driver for real-time in-memory threat detection, targeting VAD manipulation, module stomping, and anomalous executable memory regions. Evaluated against 12 controlled scenarios and detected cases missed by PE-sieve. Published as an arXiv preprint.
\href{https://arxiv.org/abs/2508.03879}{Paper}
}

\vspace{0.7mm}

\cventry{}{LogiLinux}{}{}{}{
1st place at LauzHack '25. Reverse-engineered Logitech HID++ protocols and built a Linux driver for the MX Creative Console, including a C++ backend, Tauri UI, and plugin SDK.
\href{https://github.com/logilinux/}{Project}
}

\vspace{0.7mm}

\cventry{}{Chimera: .NET Runtime Instrumentation}{}{}{}{
Reverse-engineered CLR JIT internals to study managed runtime allocation and code heap behaviour. Built a runtime instrumentation proof-of-concept using internal JIT allocation paths and documented the technique in a technical writeup.
\href{https://arjunjuneja.com/blog/2025/08/08/chimera}{Writeup}
}

\vspace{0.7mm}

\cventry{}{Source2 Engine Analysis}{}{}{}{
Analyzed Valve Source2 internals and built tooling for runtime interface enumeration, memory layout inspection, and SDK generation.
\href{https://github.com/ImArjunJ/Source2}{Project}
\href{https://arjunjuneja.com/blog/2024/03/08/source2-re}{Writeup}
}

\section{Competitions}

\begin{itemize}
    \item \textbf{Wins}: 
    \href{https://www.linkedin.com/posts/arjunjuneja_hi-just-won-pwned-7-a-ctf-ran-by-the-amazing-share-7438996490213064704-VwRX}{pwnEd 7 CTF}
    (pwn/rev, 2026),
    \href{https://www.linkedin.com/posts/arjunjuneja_hi-again-i-just-won-one-of-europes-biggest-activity-7395862662816952321-cJR1}{Junction '25}
    (\texteuro10k),
    \href{https://www.linkedin.com/posts/arjunjuneja_hi-again-yes-its-the-2nd-win-in-a-week-activity-7398432080323375104-Kj4x}{LauzHack '25}
    (EPFL), DevHouse SF, SotonHack '25, JumpStart '24.

    \item \textbf{Other placements}: ICHack '26 -- 3rd + Best Technical; RoboHack '25 -- 2nd; ICHack '25 -- 4th.

    \item \textbf{Meta Hacker Cup}: Top 6.5\% globally, placing 1465 / 22,494.
    \href{https://www.facebook.com/codingcompetitions/hacker-cup/2024/certificate/968836085258124}{Certificate}
\end{itemize}

\section{Education}

\customcventry{Sep 2024 -- Jul 2027}{BSc Computer Science}{University of Southampton}{Southampton, UK}{}{
Year 1: First Class
}

\vspace{1.5mm}

\customcventry{Sep 2022 -- Jul 2024}{A-Levels: Mathematics, Physics, Chemistry}{Feltham College}{London, UK}{}{
A*AA
}

\section{Technical Skills}

\cvitem{Languages}{C++, Rust, C\#/.NET, Python, x86/x64 Assembly, LLVM IR}
\cvitem{Systems}{Windows Kernel, Win32 API, Linux drivers, HID, CAN bus}
\cvitem{Security/RE}{IDA Pro, Ghidra, x64dbg, GDB, Wireshark, runtime analysis, binary instrumentation}
\cvitem{Engineering}{CMake, Docker, Git, CI/CD, Dear ImGui, Tauri, libusb}

\end{document}